\scanpage{99}
\mspage{17}
C'est la condition de compatibilité avec la ``structure à lacets'' --- dans le cas où $X$ est une courbe algébrique, disons (sinon, on sait bien qu'il faut introduire, dans le cas général, une ``structure à l'infini'' des groupoïdes fondamentaux \ldots). \margin{$\Vert$} Cette compatibilité \inserted{justement} s'exprime de façon particulièrement naturelle et élégante dans le contexte des groupoïdes, \unsure{$\{$} comme une condition de la nature d'une compatibilité à des morphismes de schémas (cf (39)) --- mais on \unsure{pourrait|prendrait} ici non des schémas de t.f. mais des schémas locaux épointés ou ``modèles'' épointés --- je vais revenir sur ce point de vue plus bas.

\uline{Remarque.} Revenant aux conditions a) b), il est clair d'emblée qu'elles sont indéchiffrables \inserted{(i.e. inexprimables)} en termes $\pm$ directs, comme des conditions portant directement sur les ``paramètres'' du type $g_e(u)$ (des paramètres \deleted{\unsure{comme}} \inserted{tels que} ceux désignés par $\alpha,\beta,\xi,\eta$ dans mon étude des autom. à lacets de $\pi_3$ \ldots) qui ``représentent'' un $u\in\Gamma$ (ici $\Gamma=\operatorname{Gal}\overline{\mathbf{Q}}/\mathbf{Q}$). \deleted{Pourtant} C'est pour cette raison que \unsure{le|la}
\mspage{98}
